%==============================================================
% ABSTRAK DAN ABSTRACT
% Satu halaman berisi keduanya, tanpa nomor halaman.
% Format sesuai PPKI IPB Edisi ke-4.
%==============================================================

\clearpage
\thispagestyle{empty}

% --- ABSTRAK (Bahasa Indonesia) ---
\begin{center}
{\fontsize{14pt}{17pt}\selectfont\bfseries ABSTRAK}
\end{center}

\vspace{6pt}

NAMA MAHASISWA. Judul Skripsi.
Dibimbing oleh NAMA PEMBIMBING 1 dan NAMA PEMBIMBING 2.

\vspace{12pt}

% --- ABSTRACT (Bahasa Inggris) ---
\begin{center}
{\fontsize{14pt}{17pt}\selectfont\bfseries ABSTRACT}
\end{center}

\vspace{6pt}

STUDENT NAME. Title of Thesis (skripsi).
Supervised by NAME of 1st SUPERVISOR and NAME of 2nd SUPERVISOR.

\vspace{12pt}

Narasi disusun dalam satu paragraf, isi tidak lebih dari 200 kata,
dan ditulis dalam satu halaman untuk abstrak dan abstract.
Abstrak memuat latar belakang permasalahan (tentatif), tujuan penelitian,
metode, hasil penelitian dengan penekanan pada temuan baru, dan implikasi
yang disajikan secara informatif dan faktual.
Tidak diperbolehkan mengacu pustaka, gambar, dan tabel.
Singkatan hanya dikenalkan jika masih digunakan lagi dalam bagian lain
Abstrak/Abstract.

\vspace{12pt}

\noindent
Kata kunci: ditulis dalam bahasa Indonesia, disusun berdasarkan abjad,
maksimum lima kata atau frasa

\vspace{6pt}

\noindent
\textit{Keywords}: ditulis dalam bahasa Inggris, disusun berdasarkan abjad,
maksimum lima kata atau frasa
